\section{Metodología de diseño}\label{sec:metodologia-diseno}

El diseño de la solución se estructuró en dos componentes complementarios: el diseño mediante diagramas y el diseño de las configuraciones. Esta división permitió abordar la solución desde una perspectiva arquitectónica y, posteriormente, desde una perspectiva técnica, manteniendo una relación entre la representación de los componentes y las configuraciones necesarias para su implementación.

La primera parte correspondió al diseño mediante diagramas, en el cual se representaron los componentes que conforman la solución y las relaciones existentes entre ellos. Para ello, se empleó ArchiMate como lenguaje de modelado, de acuerdo con los elementos y relaciones definidos para representar la arquitectura propuesta.

La segunda parte correspondió al diseño de las configuraciones, en la cual se especificaron los aspectos técnicos necesarios para materializar la arquitectura definida en los diagramas. Esta parte contempló las configuraciones requeridas para el funcionamiento del servicio de directorio, así como aquellas necesarias para su integración con los servicios considerados dentro del alcance.

De esta manera, ambas partes del diseño se complementaron: los diagramas permitieron definir y comunicar la estructura de la solución, mientras que el diseño de configuraciones permitió establecer los aspectos técnicos necesarios para su posterior implementación.
